% !TEX program = xelatex
% ============================================================
%  بطاقات تعليمية — Flashcards Template
%  بطاقات تعليمية قابلة للطباعة على وجهين، بشبكة 2×4 لكل صفحة
%  (سؤال في الوجه الأمامي وجواب في الوجه الخلفي) مع خطوط قص.
%  Compile with: xelatex main.tex (run twice)
% ============================================================
%  Features:
%    - A4, 11pt, article class
%    - 2x4 grid (8 cards per page), front = question, back = answer
%    - 8 sample cards (Arabic term | English term | definition)
%    - Cutting guide lines (dashed rules between cards)
%    - tcolorbox-based card cells
%    - Print double-sided (flip on long edge) then cut along guides
% ============================================================

\documentclass[11pt,a4paper]{article}

% ---- Geometry ----
\usepackage[a4paper,margin=1.5cm,top=1.5cm,bottom=1.5cm,headheight=16pt]{geometry}

% ---- Core packages ----
\usepackage{graphicx}
\usepackage{array}
\usepackage{booktabs}
\usepackage{tabularx}
\usepackage{multirow}
\usepackage{enumitem}
\usepackage{float}
\usepackage{titlesec}
\usepackage{fancyhdr}
\usepackage{setspace}
\usepackage{xcolor}
\usepackage{amsmath}
\usepackage{tcolorbox}

% ---- RTL / Arabic language support (order matters) ----
\usepackage{bidi}
\usepackage[no-math]{fontspec}
\usepackage[quiet]{polyglossia}

% ---- Fonts ----
\setmainfont[Scale=1]{NotoKufiArabic-Regular.ttf}
\newfontfamily\arabicfont[Script=Arabic,Scale=0.95]{NotoKufiArabic-Regular.ttf}
\newfontfamily\englishfont[Scale=0.95]{TeX Gyre Termes}

% ---- Languages ----
\setdefaultlanguage[calendar=gregorian,hijricorrection=1]{arabic}
\setotherlanguage[variant=british]{english}

% ---- Hyperref (load last, after polyglossia/bidi) ----
\usepackage[breaklinks,hidelinks]{hyperref}

% ---- Colors ----
\definecolor{cardframe}{HTML}{1F3A5F}
\definecolor{cardbg}{HTML}{F7F9FC}
\definecolor{ansframe}{HTML}{2E7D32}
\definecolor{ansbg}{HTML}{F1F8F1}
\definecolor{guideline}{HTML}{999999}

% ---- Line spacing ----
\setstretch{1.2}

% ---- Captions in Arabic ----
\addto\captionsarabic{%
  \renewcommand{\figurename}{الشكل}%
  \renewcommand{\tablename}{الجدول}%
}

% ---- Headers / footers ----
\pagestyle{fancy}
\fancyhf{}
\renewcommand{\headrulewidth}{0pt}
\renewcommand{\footrulewidth}{0pt}
\fancyfoot[C]{\scriptsize بطاقات تعليمية (Flashcards) --- \thepage}

% ============================================================
% Card dimensions: 2 columns x 4 rows on A4
% Usable width ~ 18cm => each card ~ 8.8cm
% Usable height ~ 25.7cm => each card ~ 6.2cm
% ============================================================
\newlength{\cardw}
\setlength{\cardw}{8.6cm}
\newlength{\cardh}
\setlength{\cardh}{5.6cm}

% ---- Front card (question side) ----
\newcommand{\cardfront}[3]{%
  % #1 = card number, #2 = Arabic term, #3 = English term
  \begin{tcolorbox}[
    colback=cardbg, colframe=cardframe,
    boxrule=0.7pt, arc=3pt,
    width=\cardw, height=\cardh,
    valign=center, halign=center,
    left=6pt, right=6pt, top=6pt, bottom=6pt,
  ]
    {\scriptsize\color{cardframe}\textbf{بطاقة #1}}\\[0.4em]
    {\Large\bfseries #2}\\[0.3em]
    {\normalsize\englishfont\itshape #3}
  \end{tcolorbox}%
}

% ---- Back card (answer side) ----
\newcommand{\cardback}[2]{%
  % #1 = card number, #2 = definition
  \begin{tcolorbox}[
    colback=ansbg, colframe=ansframe,
    boxrule=0.7pt, arc=3pt,
    width=\cardw, height=\cardh,
    valign=center, halign=center,
    left=6pt, right=6pt, top=6pt, bottom=6pt,
  ]
    {\scriptsize\color{ansframe}\textbf{الجواب #1}}\\[0.3em]
    {\small #2}
  \end{tcolorbox}%
}

% ---- Empty spacer cell (keeps grid aligned) ----
\newcommand{\emptyslot}{%
  \rule{\cardw}{\cardh}%
}

% ---- Horizontal cutting guide (dashed) ----
\newcommand{\cutguideH}{%
  \noindent\textcolor{guideline}{\rule{\linewidth}{0.4pt}}%
}

\begin{document}

% ============================================================
% FRONT PAGE — Questions (8 cards, 2x4 grid)
% ============================================================
\thispagestyle{fancy}

\begin{center}
{\large\bfseries بطاقات تعليمية --- الوجه الأمامي (Flashcards --- Front)}\\[0.2em]
{\scriptsize اطبع الوجهين ثم قص على الخطوط (Print double-sided, then cut along the lines)}
\end{center}
\vspace{0.2em}

% TODO: استبدل المحتوى أدناه ببطاقاتك الفعلية (مصطلح عربي | مصطلح إنكليزي | تعريف).
% Grid: 2 columns x 4 rows. bidi orders columns right-to-left.
\begin{center}
\setlength{\tabcolsep}{4pt}
\renewcommand{\arraystretch}{0}
\begin{tabular}{@{}c@{\hspace{6pt}}c@{}}
\cardfront{1}{الجذر}{Square Root} &
\cardfront{2}{المعادلة}{Equation} \\[3pt]
\cardfront{3}{الزاوية}{Angle} &
\cardfront{4}{النسبة}{Ratio} \\[3pt]
\cardfront{5}{الاحتمال}{Probability} &
\cardfront{6}{المتجه}{Vector} \\[3pt]
\cardfront{7}{النهاية}{Limit} &
\cardfront{8}{المشتقة}{Derivative} \\
\end{tabular}
\end{center}

% ---- Cutting guides overlay ----
\vspace{0.1em}
\begin{center}
\textcolor{guideline}{\rule{0.92\linewidth}{0.3pt}}
\end{center}

\newpage

% ============================================================
% BACK PAGE — Answers (8 cards, 2x4 grid, mirrored columns)
% Columns are reversed so answers align with questions when
% printed double-sided and flipped on the long edge.
% ============================================================
\thispagestyle{fancy}

\begin{center}
{\large\bfseries بطاقات تعليمية --- الوجه الخلفي (Flashcards --- Back)}\\[0.2em]
{\scriptsize الجواب يقابل السؤال في البطاقة ذاتها عند الطباعة على الوجهين}
\end{center}
\vspace{0.2em}

% TODO: تأكد من تطابق ترقيم الجواب مع السؤال في الوجه الأمامي.
\begin{center}
\setlength{\tabcolsep}{4pt}
\renewcommand{\arraystretch}{0}
\begin{tabular}{@{}c@{\hspace{6pt}}c@{}}
\cardback{2}{العبارة الرياضية التي تساوي بين كميتين أو أكثر وتُحل لإيجاد المجاهيل.} &
\cardback{1}{القيمة التي إذا ضُربت في نفسها أعطت العدد الأصلي، مثال: $\sqrt{9}=3$.} \\[3pt]
\cardback{4}{المقارنة بين كميتين باستخدام القسمة، ويُكتب بالصورة $a:b$ أو $a/b$.} &
\cardback{3}{انعطاف بين شعاعين يخرجان من نقطة واحدة، ويُقاس بالدرجات أو الراديان.} \\[3pt]
\cardback{6}{كمية لها مقدار واتجاه، وتُمثل بسهم في الفضاء.} &
\cardback{5}{قياس مدى احتمال وقوع حدث، ويتراوح بين $0$ و $1$.} \\[3pt]
\cardback{8}{معدل التغيّر اللحظي لدالة بالنسبة لمتغيّرها، وتُحسب عبر التفاضل.} &
\cardback{7}{القيمة التي تقترب منها الدالة عندما يقترب المتغيّر من قيمة معيّنة.} \\
\end{tabular}
\end{center}

\vspace{0.1em}
\begin{center}
\textcolor{guideline}{\rule{0.92\linewidth}{0.3pt}}
\end{center}

\end{document}
